\documentclass[11pt]{article}
\usepackage[margin=1in]{geometry}
\usepackage{amsmath,amssymb,amsthm,mathtools}
\usepackage[hidelinks]{hyperref}
\newtheorem{theorem}{Theorem}
\newtheorem{lemma}[theorem]{Lemma}
\newtheorem{corollary}[theorem]{Corollary}
\title{Explicit higher-gap adjoints and the exact repeated-root rank condition}
\author{Internal algebraic proof and exact replay; no priority claim}
\date{17 November 2026}
\begin{document}
\maketitle
The \href{CUBIC_ADJOINT_STABILITY_SFC.tex}{cubic proof} closes
$X(A+BX^2)$ at every start and positive Gamma shape. This note derives
the corresponding structure at every integer gap $d\ge2$ and identifies
the additional rank obligation. It does not prove that obligation at all
higher gaps, or close general FC/unrestricted SFC.

Put
\[
 T_{n,\gamma,d}(t)=\sum_{k=0}^n\binom nk(\gamma)_{dk}t^k,
 \quad R(z)=z^{dn}T_{n,\gamma,d}((-1)^{d+1}z^{-d}),\quad M=\gamma+dn.
                                                               \tag{1}
\]
Here $n\ge1$, $\gamma>0$, $d\ge2$. In particular $R(0)\ne0$.
With $D=D_z$ and $\delta=zD$, coefficient comparison gives
\[
 LR=0,\qquad
 L=z^{-d-1}\left[\delta\prod_{\ell=0}^{d-1}(\delta-M-\ell)
                        -z^d(\delta-dn)\right].               \tag{2}
\]
This is monic of order $d+1$, with coefficient
$a_d=d(1-M)/z$ on $D^d$. All nonzero finite points are ordinary.
A nonzero root of $R$ has multiplicity at most $d$, by its lowest
Taylor term in (2).

\section{A uniform degree-$n$ auxiliary}
\begin{theorem}\label{thm:aux}
Define $B=B_{n,\gamma,d}$ by either of the equivalent formulas
\[
 B(z)=\sum_{j=0}^n n_{\underline j}c_jz^{n-j},\qquad
 \sum_{j\ge0}c_jx^j=(1-x)^{-\gamma}
 \left(\frac{1-(1-x)^d}{dx}\right)^{-n-1}.                    \tag{3}
\]
Then $B_{0,\gamma,d}=1$ and
\[
 B_{n+1,\gamma,d}=(z+dn+\gamma+d-1)B_{n,\gamma,d}
                                  -zB_{n,\gamma,d}'.          \tag{4}
\]
For every $d$th root of unity $\omega$ the function
\[
 f_\omega(z)=e^{-\omega z}z^{1-M}B(\omega z)                   \tag{5}
\]
is a solution of the formal adjoint $L^*f=0$, on local branches away
from zero. These $d$ solutions are linearly independent.
\end{theorem}
\begin{proof}
The following coefficient change is the formal-residue version of
Lagrange inversion. Set $v=1-(1-x)^d$ in (3), noting $v=dx+O(x^2)$.
It gives
\[
 B_{n,\gamma,d}(z)=d^n n![v^n](1-v)^{-(\gamma+d-1)/d}
                       e^{z(1-(1-v)^{1/d})}.                 \tag{6}
\]
Thus its exponential generating function in $n$ is
\[
 F(t,z)=(1-dt)^{-(\gamma+d-1)/d}
                    e^{z(1-(1-dt)^{1/d})}.
\]
The identity $(1-dt)F_t=(z+\gamma+d-1)F-zF_z$ proves (4).
In particular every $B$ is monic of degree $n$.

Write $Q=1-D$. The polynomial equation needed for (5) is
\[
 \mathcal E_d B:=z(Q^{d+1}-Q)B+(M-1)(Q^d-1)B+dnB=0.        \tag{7}
\]
For completeness put
$A(x)=(1-x)((1-x)^d-1)$ and
$b(x)=(M-1)((1-x)^d-1)+dn$.
The series $C(x)$ in (3) satisfies
\[
 A C'=((n+1)A/x-A'+b)C.
\]
Taking the residue of the derivative of
$e^{zx}A(x)C(x)x^{-n-1}$ proves
$[zA(D)+b(D)]B=0$, which is (7). Every operation is on formal series
at zero; no convergence at any other singularity is used.

Equation (2) also has the divergence form
\[
 L=z^{-d}D z^{M+d}D^d z^{-M}-D+dn/z.
\]
Consequently
\[
 L^*=(-1)^{d+1}z^{-M}D^d z^{M+d}D z^{-d}+D+dn/z.
\]
Direct differentiation now gives the exact gauge identity
\[
 L^*(e^{-z}z^{1-M}B)=e^{-z}z^{-M}\mathcal E_dB.               \tag{8}
\]
Rotational covariance under $z\mapsto\omega z$, $\omega^d=1$,
proves (5), with only a nonzero branch constant for the power of $z$.
Independence follows from that of exponential polynomials with distinct
exponents: after division by $z^{1-M}$ these are
$e^{-\omega z}B(\omega z)$. For example, on a ray where a selected
exponent has uniquely largest real part its degree-$n$ leading term
cannot cancel the others; each vertex exponent has such a ray.
\end{proof}

\section{Exactly what a repeated root requires}
Let $\omega_0,\ldots,\omega_{d-1}$ be all $d$th roots of unity and put
\[
 \mathcal J_{j,k}(z)=\omega_j^k(Q^kB)(\omega_jz),
 \quad 0\le j<d,\quad 0\le k\le d-2.                       \tag{9}
\]
Using $D-1$ instead of $Q=1-D$ only changes column signs.
\begin{theorem}\label{thm:rank}
At each $z_0\ne0$, $R$ has a repeated root if and only if
$\mathcal J(z_0)$ has rank smaller than $d-1$.
Equivalently all its $(d-1)$-row minors vanish there.
Thus $T_{n,\gamma,d}$ is squarefree if and only if those minors have
no common nonzero root. This is an exact finite certificate for a
specified index and parameter, and an unproved uniform obligation at
higher gaps.
\end{theorem}
\begin{proof}
For $L=\sum_{k=0}^{d+1}a_kD^k$, $a_{d+1}=1$, the Lagrange concomitant
\[
 \mathcal C(f,R)=\sum_{k=1}^{d+1}\sum_{j=0}^{k-1}
                 (-1)^j(a_kf)^{(j)}R^{(k-1-j)}              \tag{10}
\]
has derivative $fLR-RL^*f=0$. For (5), it is an exponential times
$z^{1-M}$ times a rational function, so it vanishes identically:
otherwise a rational logarithmic derivative would tend to the nonzero
constant $\omega$ at infinity. Local powers with noninteger $M$ cause
no difficulty for this rational identity.

The $d$ concomitants form a rank-$d$ matrix on the $(d+1)$-component
jet $(R,R',\ldots,R^{(d)})$. Indeed the change from the adjoint jets
to their concomitant rows is invertible triangular, and independence
of $d$ solutions of an order-$d+1$ ordinary equation gives rank $d$
at every ordinary point. Otherwise a nonzero linear combination would
have all initial derivatives through order $d$ equal to zero.
The nullspace is therefore one-dimensional, spanned by the nonzero
jet of $R$.

Restricting the columns to $R'',\ldots,R^{(d)}$ has a nonzero nullvector
exactly when $R(z_0)=R'(z_0)=0$. Its entries are an invertible triangular
combination of $f_\omega,f_\omega',\ldots,f_\omega^{(d-2)}$.
Removing the common power gauge is another invertible triangular
column operation, and removing the exponential gauges scales rows by
nonzero factors. The resulting matrix is (9), up to column signs.
Finally $t=(-1)^{d+1}z^{-d}$ is unramified away from zero.
\end{proof}

\begin{corollary}\label{cor:det}
For $\widehat{\mathcal J}$ defined as in (9) with $0\le k\le d-1$,
\[
 \det\widehat{\mathcal J}(z)=
 \left(\prod_j\omega_j^n\right)
 \left(\prod_{i<j}(\omega_j-\omega_i)\right)R(z).              \tag{11}
\]
\end{corollary}
\begin{proof}
A standard cofactor identity can be checked directly here. For a monic
order-$d+1$ operator with coefficient $a_d$ on $D^d$, the Wronskian
of $d$ independent adjoint solutions whose concomitants with $R$ vanish
is a constant times $e^{\int a_d}R$.
To see this without an additional theorem input, write the original
companion system as $Y'=AY$ and the concomitant rows as a matrix $U$.
Then $U'=-UA$. Differentiating its alternating cofactor column $v$
gives $v'=(A-\operatorname{tr}(A)I)v$.
As $\operatorname{tr}(A)=-a_d$ and $Uv=0$, it follows that
$v$ is a constant times $e^{\int a_d}(R,\ldots,R^{(d)})^T$.
Its first component is the adjoint Wronskian, up to a constant sign.

Here $e^{\int a_d}=z^{d(1-M)}$ and $\sum_j\omega_j=0$.
Stripping the power/exponential gauges from the Wronskian leaves
$\det\widehat{\mathcal J}$ up to column signs. Thus this determinant
is a constant times $R$. Its leading coefficient is the displayed
Vandermonde product, since $B$ is monic.
\end{proof}

\section{The first higher-gap sector target and a failed stronger shortcut}
For $d=3$ put $\alpha=e^{i\pi/3}$ and define the real monic polynomial
\[
 W(z)=B(\alpha z)B(\overline\alpha z)-
 \frac{\alpha B'(\alpha z)B(\overline\alpha z)
       -\overline\alpha B'(\overline\alpha z)B(\alpha z)}{i\sqrt3}.
                                                               \tag{12}
\]
It has degree $2n$. If $R$ had a repeated root $z_0\ne0$, its three
pairwise adjoint Wronskians would vanish. Stripping the gauges gives
\[
 W(-z_0)=W(-\zeta z_0)=W(-\zeta^2 z_0)=0,
 \qquad \zeta=e^{2\pi i/3}.                                 \tag{13}
\]
Consequently zero-freeness of $W$ on the CLOSED sector
$|\arg z|\le\pi/3$ would suffice for simplicity. Every nonzero
three-point rotation orbit meets that sector. This sector assertion
is a proposed sufficient target, not a theorem established here.
For exact rational parameters, (13) is equivalently a common nonzero
root of the three polynomials $W_r(u)$ in
$W(z)=W_0(z^3)+zW_1(z^3)+z^2W_2(z^3)$.

Strict real-root interlacing of any two $W_r$ would also suffice.
The replay verifies negative simple roots and pairwise strict interlacing
for $n=\gamma=4,6,8$, by exact Sturm counts and disjoint rational
isolating intervals. This is finite evidence only; no uniform interlacing
theorem is claimed. A generalized Hurwitz-matrix route would likewise
need its own positivity proof and the appropriate boundary exclusion.

The stronger Hurwitz assertion for $W$ is false already at
$n=4,\gamma=4$ (the ordinary lower-one terminal parameter):
\[
 B=z^4+36z^3+520z^2+3560z+9720.
\]
Its $W$ is
\begin{align*}
 z^8+36z^7+740z^6+10560z^5+113800z^4
 +933920z^3+5768000z^2+24494400z+59875200.
\end{align*}
The seventh leading Hurwitz determinant is
$-169684039754781133897728000000<0$.
This disproves the Hurwitz shortcut, not squarefreeness: the three
$W_r$ have gcd one, and $T_{4,4,3}$ is squarefree by an exact gcd.

\section{A uniform initial block from the third central moment}
\begin{theorem}\label{thm:initial}
For every integer $d\ge1$, every $\gamma>0$, and $1\le n\le4$,
$T_{n,\gamma,d}$ is squarefree. Consequently $X(A+BX^d)$ has no
consecutive zero Gamma moments with terminal order at most four,
for every positive shape and every gap.
\end{theorem}
\begin{proof}
Let $Y$ have unit-rate Gamma shape $\gamma$ and put $X=Y^d$.
The reversed moment polynomial is $P_n(z)=\mathbb E(z+X)^n$.
For even $n$ it is strictly positive on the real axis. For odd $n$
its derivative $n\mathbb E(z+X)^{n-1}$ is strictly positive there.
Thus every real root is simple, and a nonreal repeated root brings its
conjugate with the same multiplicity. This already excludes repetition
for $n\le3$.

If the real monic quartic $P_4$ had a repeated root, it would have to be
$(z^2+sz+t)^2$. With $\mu_j=\mathbb EX^j$, comparison of its first
three coefficients gives $s=2\mu_1$, $t=3\mu_2-2\mu_1^2$ and
$\mu_3=3\mu_1\mu_2-2\mu_1^3$. This would give zero third central moment.
But put $r(x)=(x)_d$, $a=r(\gamma)$, $b=r(\gamma+d)$,
$c=r(\gamma+2d)$. The polynomial $r$ is strictly increasing and convex
on the positive real axis: its coefficients are nonnegative, with
$r'>0$ and $r''\ge0$. Hence $b>a>0$ and $c\ge2b-a$, giving
\[
 \mathbb E(X-\mathbb EX)^3
   =abc-3a^2b+2a^3\ge2a(b-a)^2>0.
\]
This contradiction proves the quartic case. Reversal preserves
multiplicity since the constant and leading coefficients are nonzero.
The consecutive-moment bridge below gives the stated SFC initial block.
No assertion about terminal orders beyond four follows from this argument.
\end{proof}

\section{SFC bridge and remaining scope}
For every gap $d\ge1$, every positive shape $\beta$, and terminal
$n\ge1$, put $\gamma=n+\beta-1$ and $\theta=tD_t$. The identities
\[
 T_{n,\gamma+1,d}=T_{n,\gamma,d}+(d/\gamma)\theta T_{n,\gamma,d},
 \qquad T_{n-1,\gamma,d}=T_{n,\gamma,d}-(1/n)\theta T_{n,\gamma,d}
\]
identify consecutive Gamma moments of $X(A+BX^d)$ with the
multiple-root obstruction. Their determinant is
$-(\gamma+dn)/(\gamma n)\ne0$.
The rank theorem therefore supplies the precise uniform task for this
whole class. It does not supply its solution at $d\ge3$.
The exact checker replays the algebra and includes nonpositive/complex
parameter controls where a repeated root genuinely occurs, so the rank
certificate is tested on both outcomes. Those controls are outside
positive Gamma scope and are not SFC counterexamples.
All novelty, independent/formal/external assurance flags remain false.
\end{document}
